\section{Failure and null outcomes as persistent knowledge}

Self-ORION stores failed, blocked, partial and \texttt{CANNOT\_CHECK} executions as addressable experience rather than disposable execution noise. A failure episode distinguishes the observed mode, propagated effect, task/variation identity, competing causes, detection evidence, recovery actions, residuals and falsifiers. Negative outcomes remain in history after later success.

This makes failure a potential research input, but not an automatic method lesson. Let $e_j$ denote an immutable episode and $\phi(e_j)$ its observed failure signature. Repeated signatures can justify a recurrence hypothesis,
\[
\phi(e_1)\approx\phi(e_2)\approx\cdots\approx\phi(e_n),
\]
without justifying the causal statement that one mechanism caused every failure. Recurrence answers ``does something structurally similar persist?''; attribution answers ``which intervention should change the outcome?''

A valid development record therefore keeps failures, nulls and harmful interventions visible. Otherwise an adaptive system can manufacture an apparently improving history by retaining only successful challengers. Self-ORION treats negative-history completeness as an assurance coordinate rather than an optional debugging convenience.

\subsection{Boundary with Paper I}

Paper I owns the mechanic-cell representation that can expose a missing or failed coordinate. Paper V begins when execution evidence is interpreted as possible method evidence. The transition from an open mechanic coordinate to a method change is deliberately non-automatic: a missing specification, a provider outage, a data defect and a genuine operator defect require different responses.
